\documentclass[../../../include/open-logic-section]{subfiles}

\begin{document}

\olfileid{his}{set}{infinitesimals}
\olsection{انفينټېسيمالونه او مشتق اخيستل}

نيوټن او لايبنيڅ د 17مې پېړۍ په پای کښې، له يو بل څخه په خپلواکه توګه، حسابان وموندل۔ د حسابان يوه په ځانګړې توګه مهمه کارونه \emph{مشتق اخيستل} وو۔ په لنډه توګه، د مشتق اخيستلو موخه دا ده چې په يوه نقطه کښې د يوې تابع د «بدلون د کچې»، يا ميلان، يو مفهوم راکړي۔

د دې نظر د څرګندولو يوه روښانه لاره دا ده۔ تابع
$f(x) = \nicefrac{x^2}{4} + \nicefrac{1}{2}$ په پام کښې ونيسئ، چې لاندې په تور رنګ ښودل شوې ده:
\begin{center}
	\begin{tikzpicture}[scale=1]
	\draw[->, gray] (-1,0) -- (4.25,0) node[right] {$x$};
	\draw[->, gray] (0,-0.25) -- (0,5.25) node[above] {$f(x)$};

	\foreach \x/\xtext in {1/1, 2/2, 3/3, 4/4}
	\draw[shift={(\x,0)}, gray] (0pt,2pt) -- (0pt,-2pt) node[below] {$\xtext$};

	\foreach \y/\ytext in {1/1, 2/2, 3/3, 4/4, 5/5}
	\draw[shift={(0,\y)}, gray] (2pt,0pt) -- (-2pt,0pt) node[left] {$\ytext$};

	\draw[oldiagcolorC] (0.5, 9/16) -- (3.5, 57/16) -- (3.5, 9/16)--cycle;
	\draw[oldiagcolorD] (.5, 9/16) -- (2.5, 33/16) -- (2.5, 9/16) -- cycle;
	\draw[oldiagcolorE] (.5, 9/16) -- (1.5, 17/16) -- (1.5, 9/16) -- cycle;
	\draw[thick] (-1,.75) parabola bend (0,0.5) (4,4.5);
	\end{tikzpicture}
\end{center}
فرض کړئ چې غواړو په $c =
\nicefrac{1}{2}$ کښې د تابع ميلان پيدا کړو۔ پيل په داسې مثلث کوو چې وتر يې په هغه نقطه کښې د ميلان تقريب راکوي؛ کېدای شي دا هماغه پورتنے سور مثلث وي۔ که $\beta$ د مثلث د قاعدې اوږدوالے وي، نو لوړوالے يې
$f(\nicefrac{1}{2}+\beta) - f(\nicefrac{1}{2})$ دے؛ ځکه نو د وتر ميلان دا دے:
\[
\frac{f(\nicefrac{1}{2}+\beta) - f(\nicefrac{1}{2})}{\beta}.
\]
نو زموږ د !!{colorC} مثلث ميلان، چې د قاعدې اوږدوالے يې~$3$ دے، عين~$1$ دے۔ د يو کوچني مثلث وتر، يعنې د !!{colorD} مثلث چې د قاعدې اوږدوالے يې~$2$ دے، غوره تقريب راکوي؛ ميلان يې $\nicefrac{3}{4}$ دے۔ تر دې هم کوچنے !!{colorE} مثلث، چې د قاعدې اوږدوالے يې~$1$ دے، لا غوره تقريب راکوي؛ ميلان يې $\nicefrac{1}{2}$ دے۔

څومره چې مثلثونه کوچني کېږي، هومره غوره تقريبونه راکوي۔ نو داسې څه ويلے شو: د هغه مثلث وتر چې د قاعدې اوږدوالے يې \emph{انفينټېسيمال}، يعنې بې‌نهايت کوچنے، وي، په خپله د $c =
\nicefrac{1}{2}$ ميلان راکوي۔ په دې توګه به د تابع $f$ د (لومړي) مشتق لپاره، په نقطه $c$ کښې، دا فارمول ترلاسه کړو:
\[
{f'}(c) = \frac{f(c+\beta) - f(c)}{\beta} \text{ چې $\beta$ انفينټېسيمال دے۔}
\]
او په تقريباً همدغه معنا نيوټن او لايبنيڅ دا خبره کړې وه۔

خو چې هغوي دا خبره کړې ده، نو بايد ترې وپوښتو: \emph{انفينټېسيمال} څه دے؟ يوه جدي دوه‌لارې مسئله پيدا کېږي۔ که $\beta = 0$ وي، نو $f'$ سم نۀ دے تعريف شوے، ځکه په $0$ وېش پکې راځي۔ خو که $\beta >
0$ وي، نو يوازې د ميلان يو \emph{تقريب} ترلاسه کوو، خپله ميلان نۀ ترلاسه کوو۔

دا د پخواني وخت په اړه د وروستي وخت يوه بې‌ځايه اندېښنه نۀ ده۔ برکلي د نيوټن پر پيروانو داسې نيوکه کوي:
\begin{quote}
زه منم چې نښې د هر څه يا د هېڅ شي لپاره ټاکل کېدای شي؛ نو په اصلي نښه‌ليکنه $c + \beta$ کښې $\beta$ هم يو زياتوالے او هم هېڅ ښودلے شو۔ خو په دې دواړو کښې چې هره معنا ورته وټاکئ، بايد د هماغې معنا سره سم يو شان استدلال وکړئ، او په دوه‌ګونې معنا استدلال مخ ته يونۀسئ؛ دا کار به ښکاره مغالطه وي۔ (\citealt[\S{}XIII]{Berkeley1734}، متغيرونه د پورتني متن سره د سمون لپاره بدل شوي دي۔)
\end{quote}
له برکلي څخه د انفينټېسيمال حسابان د دفاع لپاره څوک ويلے شي چې د «انفينټېسيمالونو» خبرې يوازې مجازي دي۔ ويلے شي چې تر څو يو \emph{په رښتيا کوچنے} مثلث واخلو، د مماس \emph{کافي ښه} تقريب به ترلاسه کړو۔ برکلي دې ته هم ځواب درلود: دا که د انجينيرۍ لپاره کافي ښه وي، خو د رياضياتو \emph{حيثيت} کمزورے کوي، ځکه
\begin{quote}
موږ ته ويل کېږي چې \emph{in rebus mathematicis errores qu\`{a}m minimi
non sunt contemnendi}. [په رياضياتو کښې تر ټولو وړې تېروتنې هم بايد له پامه ونۀ غورځول شي۔] \citep[\S{}IX]{Berkeley1734}
\end{quote}
په مايله بڼه ليکل شوے عبارت د نيوټن د خپل کتاب \emph{د منحني‌ګانو مساحت موندل} (1704) تقريباً لفظ په لفظ نقل دے۔

د برکلي فلسفي اعتراضونه ډېر ژور او تېز دي۔ سره له دې، حسابان بېخي بريالے کار و، او رياضيپوهانو له تېروتنې سره د مخ کېدو پرته د هغه کارول جاري وساتل۔

\end{document}
